\section{ROM 放不下不断增长的内核}

上一章结束时，CPU 已经能从 128 KiB ROM 启动并向 COM1 输出文字。继续把
内核代码塞进 ROM 很快会遇到两个问题：空间有限，更新也不方便。磁盘适合保存
较大的程序，但 CPU 不能直接执行“磁盘里的地址”。代码必须先被读入 RAM。

我们先把目标缩小到 Stage 1。磁盘的 LBA 0 保存描述符，负载从 LBA 1 开始。
ROM 检查描述符和校验值，把 Stage 1 复制到物理地址 \code{0x8000}，然后
跳过去执行。成功时串口依次出现：

\begin{lstlisting}
[OS][FIRMWARE] STAGE1_HEADER_VALID
[OS][FIRMWARE] STAGE1_LOADED
[OS][STAGE1] A20_READY
[OS][STAGE1] ENTERED
\end{lstlisting}

后两行由刚刚读入 RAM 的代码输出。看到它们以后，我们才知道磁盘端口、校验、
目标内存、控制转移和 A20 地址实验都已经走通。下一步是让这段 16 位代码进入
64 位模式，再装入 ELF64 内核。

\topic{先手算一次 1 MiB 地址}
实模式用“段值左移四位，再加偏移”的方法形成线性地址。取
\code{FFFF:0010}：
\[
  \mathtt{0xFFFF}\times16+\mathtt{0x0010}
  =\mathtt{0x100000}.
\]
8086 只有 20 根地址线，最高只能表达 \code{0xFFFFF}。多出来的第 21 位被
丢掉以后，结果回到 \code{0x00000}。一些旧软件依赖这种回绕，后来能够访问
更大地址空间的 PC 因而保留了 A20 兼容门。

Stage 1 需要把 Kernel 放到 1 MiB 以上，所以先通过端口 \code{0x92} 请求
打开 A20。写端口还不够：代码保存两个测试地址的原值，写入不同字节，检查它们
是否仍然互不影响，最后恢复原值。串口输出 \code{A20_READY} 时，表示这次
地址别名实验已经通过。

\topic{为什么 64 位 CPU 仍从实模式开始}
8086 的实模式直接用段值参与地址计算。80286 增加保护模式以后，段寄存器改为
选择 GDT 中的描述符；80386 又加入 32 位寄存器和分页。AMD64 在这些状态
之上增加长模式，同时保留旧入口，使新处理器仍能启动旧软件。

这里没有“64 位 CPU 先以 32 位运行，再退回实模式”的过程。RESET 直接建立
实模式状态。Stage 1 要从这个起点向前经过两个远跳转：第一次设置
CR0.PE，并装入 D=1 的 32 位代码段；第二次在页表和 IA32\_EFER.LME
准备好以后，装入 L=1 的 64 位代码段。

中间还有一个容易漏掉的状态。设置 LME、PAE 和 CR0.PG 后，
IA32\_EFER.LMA 变为 1，IA-32e 模式已经激活；但当前 CS 仍指向 L=0、
D=1 的代码段，所以 CPU 暂时按 32 位兼容子模式继续解码。只有随后的 far
jump 装入 L=1 的代码段，目标位置才按 64 位方式解码。

只修改汇编文件中的 \code{BITS 32} 或 \code{BITS 64} 不会改变 CPU 模式，
那只是告诉 NASM 应该生成哪一种指令编码。

\begin{center}
\begin{tikzpicture}[os diagram]
  \node[os state] (real) {实模式\\默认 16};
  \node[os block,right=of real] (protected) {保护模式\\CS.D=1\\默认 32};
  \node[os block,below=of protected] (compat) {IA-32e 已激活\\CS.L=0\\兼容子模式};
  \node[os state,left=of compat] (long) {CS.L=1\\64 位子模式};
  \draw[os arrow] (real) -- (protected);
  \draw[os arrow] (protected) -- node[right]{PAE、LME、PG} (compat);
  \draw[os arrow] (compat) -- node[below]{far jump} (long);
\end{tikzpicture}
\end{center}

\begin{osgridlongtable}{L{0.20\linewidth}L{0.22\linewidth}L{0.25\linewidth}L{0.23\linewidth}}
执行区间 & 决定状态的关键位 & 默认解码宽度 & 下一次变化由谁触发 \\ \hline
\endfirsthead
执行区间 & 决定状态的关键位 & 默认解码宽度 & 下一次变化由谁触发 \\ \hline
\endhead
Stage 1 前半段 & PE=0、PG=0 & 16 位 & 写 CR0.PE \\ \hline
32 位保护模式 & PE=1、PG=0、CS.D=1 & 32 位 & 设置 PAE、LME、PG \\ \hline
IA-32e 兼容子模式 & LMA=1、CS.L=0、CS.D=1 & 32 位 & far jump 装入 L=1 的 CS \\ \hline
64 位子模式 & LMA=1、CS.L=1、CS.D=0 & 64 位 & 已到 64 位入口 \\ \hline
\end{osgridlongtable}

\topic{“保护模式”保护的是什么}
\term{保护模式}{protected mode} 得名于硬件开始检查内存段、控制转移和
特权边界。实模式中的段值直接参与地址加法；保护模式中的段寄存器保存
\term{选择子}{selector}，选择子指向 GDT 或 LDT 中的描述符。描述符说明
段基址、界限、代码或数据类型、可读写性、是否存在以及允许哪个特权级访问。

“保护”不表示把 CR0.PE 置一后机器立刻安全。软件必须先建立正确的描述符表，
为内核和用户代码选择特权级，安装异常入口，并在需要时建立带权限的页表。
如果所有段覆盖整个地址空间、所有页都可写、所有程序都运行在 ring 0，CPU
虽然处于保护模式，程序之间仍没有实际隔离。

80286 首先把保护模式带入 x86，但它仍以 16 位代码为主。它把物理地址扩到
24 位，可访问 16 MiB，并引入描述符表、四个特权环、任务状态与保护式中断。
早期 286 退出保护模式的正式路径十分笨重，常需借助复位；这也说明保护模式
最初被设想为操作系统长期驻留的环境，不是每个函数临时打开的开关。

\topic{80386 怎样在不抛弃 16 位软件的前提下扩展到 32 位}
80386 没有重新设计一套互不兼容的寄存器和指令。它在原状态外面继续加宽，
再让代码段属性和指令前缀选择本条代码使用旧宽度还是新宽度：

\begin{osgridlongtable}{L{0.24\linewidth}L{0.30\linewidth}L{0.36\linewidth}}
被扩展的部分 & 80386 的做法 & 为什么旧代码仍能运行 \\ \hline
\endfirsthead
被扩展的部分 & 80386 的做法 & 为什么旧代码仍能运行 \\ \hline
\endhead
通用寄存器 & AX 外扩成 EAX，其他七个寄存器同样外扩 &
旧操作码仍可只访问低 16 位 AX \\ \hline
指令指针与标志 & IP/FLAGS 外扩成 EIP/EFLAGS &
16 位代码仍按低 16 位控制流和旧标志语义运行 \\ \hline
默认操作数和地址 & 代码段描述符增加 D 位 &
D=0 默认 16 位，D=1 默认 32 位 \\ \hline
单条指令切换宽度 & 增加 \code{66} 与 \code{67} 前缀 &
16 位段可临时用 32 位，32 位段也可临时用 16 位 \\ \hline
有效地址 & 加入 32 位 base、index、scale、displacement 和 SIB &
原 16 位 BX/BP/SI/DI 编码保持原义 \\ \hline
段基址与界限 & 描述符支持 32 位基址、20 位 limit 和 G 粒度 &
16 位段仍可保持 64 KiB，32 位平坦段可覆盖 4 GiB \\ \hline
物理地址 & 386DX 提供 32 位物理地址空间 &
分页关闭时可直接到 4 GiB；旧机器地址仍落在低端 \\ \hline
虚拟内存 & 增加 CR3 和两级 4 KiB 分页 &
分页由 CR0.PG 单独开启，不强迫旧保护模式代码使用 \\ \hline
额外段寄存器 & 增加 FS 与 GS &
CS、DS、SS、ES 的旧编码与用途不变 \\ \hline
旧 8086 程序 & 增加 virtual-8086 环境 &
保护模式操作系统可把旧程序放在受分页控制的任务中运行 \\ \hline
\end{osgridlongtable}

\begin{center}
\begin{tikzpicture}[os diagram]
  \node[os hardware,minimum width=22mm] (al) {AL\\8 位};
  \node[os state,right=of al,minimum width=25mm] (ax) {AX\\低 16 位};
  \node[os block,right=of ax,minimum width=28mm] (eax) {EAX\\32 位};
  \node[os state,below=of ax,minimum width=25mm] (ah) {AH\\bit 15:8};
  \draw[os arrow] (al) -- (ax);
  \draw[os arrow] (ah) -- (ax);
  \draw[os arrow] (ax) -- (eax);
\end{tikzpicture}
\end{center}
\diagramnote{EAX 不是与 AX 平行的第二个寄存器。写 AX 改变 EAX 的低
16 位而保留高 16 位；写 AL 或 AH 只改变相应 byte。}

这种扩展方式解释了为什么本项目能在实模式读写 EAX，也解释了进入 D=1 的
32 位代码段后，原来的 16 位形式仍可通过 \code{66} 使用。处理器模式、
代码段默认值与单条指令宽度始终是三层状态。

\topic{选择子怎样找到一个受检查的段}
保护模式的 16 位选择子不再表示 \(segment\ll4\)。其低两位是请求特权级
RPL，bit 2 是表指示 TI，其余 13 位是描述符索引：
\[
  selector=(index\ll3)\;|\;(TI\ll2)\;|\;RPL.
\]
项目的代码选择子 \code{0x0008} 拆开以后：
\[
  index=1,\qquad TI=0,\qquad RPL=0.
\]
所以它选择 GDT 第 1 项，并以 ring 0 请求访问。数据选择子
\code{0x0010} 的 index 为 2。

\begin{pdfdiagram}
  \node[os state] (selector) {段选择子\\index, TI, RPL};
  \node[os block,right=of selector] (table) {GDT 或 LDT\\8-byte 描述符};
  \node[os state,right=of table] (cache) {隐藏段缓存\\base, limit, attributes};
  \node[os block,below=of cache] (offset) {EIP/ESP/有效偏移};
  \node[os hardware,left=of offset] (linear) {线性地址\\base+offset};
  \node[os block,left=of linear] (paging) {可选分页\\linear \(\to\) physical};
  \draw[os arrow] (selector) -- (table);
  \draw[os arrow] (table) -- (cache);
  \draw[os arrow] (cache) -- (linear);
  \draw[os arrow] (offset) -- (linear);
  \draw[os arrow] (linear) -- (paging);
\end{pdfdiagram}
\diagramnote{加载段寄存器时，CPU 检查描述符并缓存 base、limit 和属性。
以后每次访问使用隐藏缓存，不会每条指令重新读取 GDT。}

8-byte 代码/数据描述符包含 32 位 base、20 位原始 limit、类型、S、DPL、
P、AVL、L、D/B 和 G。G=0 时 limit 按 byte 计；G=1 时：
\[
  effective\_limit=(raw\_limit\ll12)\;|\;0xFFF.
\]
本项目 32 位代码描述符为 \code{0x00CF9A000000FFFF}。其中 base=0，
raw limit=\code{0xFFFFF}，G=1，所以有效界限为 \code{0xFFFFFFFF}；
P=1、DPL=0、S=1、type 表示可读可执行代码，L=0、D=1。它建立的是
base 为零、覆盖 4 GiB、默认 32 位的 ring 0 代码段。

\topic{32 位保护模式的地址到底有多少位}
在本项目的平坦段模型中，CS、DS、ES 和 SS 对应描述符的 base 都为零，limit
覆盖 \code{0x00000000..0xFFFFFFFF}。因此 32 位有效偏移与线性地址数值
相同，但 CPU 仍完成描述符类型、界限和权限检查。

\begin{osgridlongtable}{L{0.24\linewidth}L{0.18\linewidth}L{0.48\linewidth}}
量 & 32 位基线 & 含义 \\ \hline
\endfirsthead
量 & 32 位基线 & 含义 \\ \hline
\endhead
段选择子 & 16 位 &
包含 13 位索引、TI 与 RPL，不直接成为地址高位 \\ \hline
段描述符 base & 32 位 &
加到有效偏移上，项目平坦段把它设为零 \\ \hline
默认有效地址 & 32 位 &
由 EAX--EDI、SIB、位移或 EIP/ESP 等形成 \\ \hline
线性地址 & 32 位 &
范围为 \code{0x00000000..0xFFFFFFFF}，共 4 GiB \\ \hline
物理地址 & 386DX 为 32 位 &
分页关闭时等于线性地址；后来 PAE 处理器可实现更多物理位 \\ \hline
页内偏移 & 12 位 &
传统 4 KiB 页的低位，不参加页表索引 \\ \hline
页目录/页表索引 & 各 10 位 &
传统两级分页用 \(10+10+12=32\) 位覆盖 4 GiB \\ \hline
默认操作数 & 8 或 32 位 &
D=1 时普通整数形式默认使用 32 位；\code{66} 可临时选择 16 位 \\ \hline
\end{osgridlongtable}

\begin{workedexample}
项目数据选择子为 \code{0x0010}，描述符 base=0、limit=
\code{0xFFFFFFFF}。执行 \code{mov eax, [ebx+esi*4+8]}，若
\code{EBX=0x00100000}、\code{ESI=3}：
\[
  offset=0x00100000+3\times4+8=0x00100014.
\]
DS base 为零，线性地址仍是 \code{0x00100014}。分页尚未开启时，物理地址
也等于 \code{0x00100014}；打开分页以后，同一个线性地址还要由 CR3 指向的
页表翻译。
\end{workedexample}

\topic{特权环、描述符和分页怎样形成保护}
x86 定义 ring 0 到 ring 3 四个特权级，数字越小权限越高。CPL 来自当前
CS 的低两位，描述符携带 DPL，选择子还携带 RPL。不同访问类型使用不同的
比较规则，因此不能把所有检查简化成“数字小就能访问”。

\begin{osgridlongtable}{L{0.20\linewidth}L{0.27\linewidth}L{0.43\linewidth}}
检查位置 & 主要字段 & 能阻止什么 \\ \hline
\endfirsthead
检查位置 & 主要字段 & 能阻止什么 \\ \hline
\endhead
代码/数据描述符 & type、DPL、P、limit &
向数据段取指、向只读段写入、越界或访问不存在的段 \\ \hline
控制转移 & CPL、RPL、DPL、门描述符 &
低权限代码直接跳入高权限入口或伪造返回状态 \\ \hline
页表 & P、RW、US，后续还有 NX &
访问未映射页、用户访问 supervisor 页或违反读写权限 \\ \hline
I/O 权限 & IOPL 与 TSS I/O bitmap &
低权限程序直接操作设备端口 \\ \hline
系统指令 & 当前模式与 CPL &
用户程序改 CR3、关闭保护或重写描述符表 \\ \hline
\end{osgridlongtable}

分页在保护模式中是可选的。PE=1、PG=0 时，描述符仍然检查访问，线性地址
直接成为物理地址；PE=1、PG=1 时，页表再做一层映射和权限检查。项目先进入
无分页的 32 位保护模式，用这个较简单的环境构造页表，之后才打开 PG。

\topic{32 位保护模式能执行哪些指令}
进入 D=1 的代码段以后，原有整数、控制流、栈、字符串和端口指令仍然存在，
但默认寄存器与偏移扩成 EAX、EIP、ESP 和 32 位有效地址。

\begin{osgridlongtable}{L{0.22\linewidth}L{0.31\linewidth}L{0.37\linewidth}}
类别 & 典型指令 & 相比实模式新增或改变的重点 \\ \hline
\endfirsthead
类别 & 典型指令 & 相比实模式新增或改变的重点 \\ \hline
\endhead
普通计算 & MOV、ADD、SUB、IMUL、DIV、AND、OR、XOR、SHIFT &
32 位操作数成为默认，16 位形式使用 \code{66} \\ \hline
32 位寻址 & LEA、所有带 ModR/M/SIB 的内存形式 &
基址加变址乘比例，再加位移，可覆盖 4 GiB 段内偏移 \\ \hline
过程与栈 & CALL、RET、PUSH、POP、ENTER、LEAVE &
默认保存 EIP，ESP 每次按操作数宽度移动 \\ \hline
保护式中断 & INT、IRET、32 位 IDT gate &
门描述符检查特权，可切换到 TSS 提供的高权限栈 \\ \hline
描述符管理 & LGDT、LIDT、LLDT、LTR、LAR、LSL、VERR、VERW &
装入或检查全局、局部和任务相关描述符状态 \\ \hline
分页与控制 & MOV CRn、INVLPG、CLTS &
选择页表根、使局部翻译失效或管理浮点任务状态 \\ \hline
系统调用入口 & call gate、INT；后续 CPU 还可提供 SYSENTER/SYSEXIT &
从 ring 3 进入 ring 0 时由硬件检查入口和栈状态 \\ \hline
原子与同步 & LOCK、XCHG、CMPXCHG、XADD、CMPXCHG8B &
为多处理器共享内存提供原子读改写基础 \\ \hline
\end{osgridlongtable}

具体指令仍受处理器代际和 CPUID 约束。保护模式不是“所有指令都可执行”的
万能状态；VMX、SSE、AVX、MSR 等各有控制位、特权与功能检查。相反，一些
系统指令在 ring 3 会产生 \(\#GP\)，正是保护模式发挥作用的表现。

\topic{项目怎样从 16 位入口进入 32 位保护模式}
下面这段正式代码横跨两个默认解码环境：

\lstinputlisting[
  language={[x86masm]Assembler},
  firstline=65,
  lastline=82,
  caption={CR0.PE 与 far jump 分别改变什么}
]{../../../../source/boot/stage1/src/entry.asm}

\begin{osgridlongtable}{L{0.19\linewidth}L{0.26\linewidth}L{0.45\linewidth}}
步骤 & 执行后的关键状态 & 为什么下一步依赖它 \\ \hline
\endfirsthead
步骤 & 执行后的关键状态 & 为什么下一步依赖它 \\ \hline
\endhead
\code{lgdt [memory]} & GDTR 指向项目 GDT，仍是实模式 &
PE 置一后，far jump 的选择子需要从这里取得描述符 \\ \hline
\code{mov eax, cr0} & EAX 保存旧 CR0 &
必须保留未打算改变的控制位 \\ \hline
\code{or eax, 1} & 只在 EAX 候选值中设置 PE &
CPU 仍在实模式，可以在写回前检查计算结果 \\ \hline
\code{mov cr0, eax} & CR0.PE=1 &
保护模式规则开始生效，但 CS 仍保留旧隐藏状态 \\ \hline
far \code{jmp 0x0008:target} & CS 选择 GDT 第 1 项，D=1 &
刷新取指路径，目标开始按 32 位默认宽度解码 \\ \hline
\code{mov ds/es/ss, 0x0010} & 数据和栈段采用 GDT 第 2 项 &
后续内存与栈访问不再依赖实模式遗留段缓存 \\ \hline
\code{mov esp, 0x7000} & 32 位栈顶确定 &
CALL 和异常才有已知的保存位置 \\ \hline
\end{osgridlongtable}

源码中的 \code{[bits 32]} 不会修改 CPU。它告诉 NASM：“下面标签将由
D=1 的代码段执行，请按 32 位默认值编码。”若 far jump 选中的描述符仍是
D=0，CPU 会用 16 位默认值解释同一串 byte，汇编器意图与机器状态便分离。

\topic{从 32 位扩展到 64 位时，哪些部分继续兼容}
AMD64 再次采用“保留旧编码，增加可选择的宽度和寄存器”的办法。AMD 手册把
新环境称为 \term{长模式}{long mode}，并把它分成 64 位子模式和兼容子模式。
所以“长模式”不是“每条指令都操作 64 位”的同义词，而是容纳新 64 位执行
环境和旧 16/32 位程序的上层状态。

\begin{osgridlongtable}{L{0.23\linewidth}L{0.31\linewidth}L{0.36\linewidth}}
扩展对象 & AMD64 的做法 & 兼容策略 \\ \hline
\endfirsthead
扩展对象 & AMD64 的做法 & 兼容策略 \\ \hline
\endhead
通用寄存器 & EAX 外扩为 RAX，并增加 R8--R15 &
旧名称继续选择低 8/16/32 位 \\ \hline
指令指针和栈 & EIP/ESP 外扩为 RIP/RSP &
兼容子模式仍按旧代码段属性执行 16/32 位程序 \\ \hline
指令编码 & 增加 REX 前缀的 W、R、X、B 位 &
旧操作码主体保留，64 位子模式才把 \code{0x40..0x4F} 解释为 REX \\ \hline
地址计算 & 默认地址宽度扩为 64 位，并增加 RIP-relative &
\code{67} 可选择 32 位地址形式，旧 SIB 结构继续使用 \\ \hline
分段 & CS/DS/ES/SS 的 base 在 64 位子模式中按零处理 &
FS/GS base 保留，兼容子模式仍使用旧分段语义 \\ \hline
分页 & PAE 表项扩成四级长模式页表 &
兼容子模式和 64 位子模式共享同一 IA-32e 页表体系 \\ \hline
控制转移与中断 & 64 位 IDT gate、RIP、RSP、TSS/IST &
旧应用可在兼容子模式运行，由 64 位内核处理事件 \\ \hline
\end{osgridlongtable}

在旧模式中，单字节 \code{0x40..0x4F} 曾编码 INC/DEC 寄存器。64 位子模式
把这段编码空间解释为 REX：
\[
  \underbrace{0100}_{REX\ 标识}\;W\;R\;X\;B.
\]
W 选择 64 位操作数；R、X、B 分别扩展 ModR/M 的 reg、SIB index 和
base/rm 字段，使原来的 3-bit 寄存器编号能够选择 R8--R15。旧 INC/DEC
仍可用通用的 ModR/M 编码表达，因此没有丢失运算能力。

\begin{workedexample}
\code{89 D8} 在 64 位子模式中表示 \code{mov eax, ebx}。加上
\code{REX.W=0x48} 后，\code{48 89 D8} 表示
\code{mov rax, rbx}。ModR/M 的源、目的关系没有重写，REX.W 只把操作数
宽度提升为 64 位。
\end{workedexample}

\topic{64 位模式中的“默认 64 位”只对了一半}
64 位子模式的默认地址宽度是 64 位，但多数普通整数指令的默认操作数宽度仍是
32 位。\code{REX.W} 通常选择 64 位操作数；\code{66} 可选择 16 位。
byte 指令继续操作 8 位。

\begin{osgridlongtable}{L{0.24\linewidth}L{0.20\linewidth}L{0.46\linewidth}}
项目 & 64 位子模式的规则 & 例子 \\ \hline
\endfirsthead
项目 & 64 位子模式的规则 & 例子 \\ \hline
\endhead
默认地址宽度 & 64 位 & \code{[rbx+rsi*4+8]} \\ \hline
\code{67} 地址覆盖 & 32 位，结果按规则扩展到线性地址 &
\code{[ebx+esi*4+8]} \\ \hline
普通整数默认操作数 & 32 位 & \code{mov eax, ebx} \\ \hline
\code{REX.W} 操作数 & 64 位 & \code{mov rax, rbx} \\ \hline
\code{66} 操作数覆盖 & 通常选择 16 位 & \code{mov ax, bx} \\ \hline
近 CALL/RET、PUSH/POP & 栈项通常为 64 位 &
RSP 每压入一个返回 RIP 减 8 \\ \hline
写 32 位通用寄存器 & 结果零扩展到完整 64 位寄存器 &
\code{mov eax, 1} 令 RAX=\code{0x0000000000000001} \\ \hline
写 16/8 位子寄存器 & 保留其余高位 &
\code{mov ax, 1} 只替换 RAX 低 16 位 \\ \hline
\end{osgridlongtable}

写 EAX 自动清 RAX 高 32 位，是 AMD64 新增的重要规则。它避免普通 32 位
结果依赖旧高位，也让 CPU 更容易消除部分寄存器依赖。写 AX、AL 或 AH 没有
同样的清高位效果。

使用任何 REX 前缀时，传统高 byte 寄存器 AH、BH、CH、DH 不能在该指令中
编码；相应编码位置用于 SPL、BPL、SIL、DIL。这是扩展寄存器编码必须付出的
兼容代价，汇编器报错时不能把它误解为寄存器本身消失。

\topic{64 位虚拟地址和物理地址都没有简单等于 64 位}
长模式要求分页开启，不存在本项目可使用的“关闭分页的 64 位子模式”。CPU
先检查线性地址是否 canonical，再用 CR3 指向的长模式页表翻译。

本项目使用四级页表，对应 48 位 canonical 虚拟地址。bit 63:48 必须全部
复制 bit 47：

\begin{osgridlongtable}{L{0.30\linewidth}L{0.19\linewidth}L{0.41\linewidth}}
地址 & 是否 canonical & 原因 \\ \hline
\endfirsthead
地址 & 是否 canonical & 原因 \\ \hline
\endhead
\code{0x00007FFFFFFFF000} & 是 &
bit 47=0，高 16 位全零 \\ \hline
\code{0xFFFF800000001000} & 是 &
bit 47=1，高 16 位全一 \\ \hline
\code{0x0000800000000000} & 否 &
bit 47=1，但高 16 位没有符号扩展为全一 \\ \hline
\code{0xFFFF7FFFFFFFFFFF} & 否 &
bit 47=0，但高 16 位不是全零 \\ \hline
\end{osgridlongtable}

四级索引使用 bit 47:39、38:30、29:21、20:12，低 12 位是 4 KiB
页内偏移。支持 LA57 的处理器还可采用五级、57 位 canonical 地址，但项目
当前没有开启 CR4.LA57。

物理地址宽度由具体 CPU 的 CPUID 报告，常见实现少于 64 位。页表项虽然是
64 位，也只有其中一段 bit 保存物理页框号，其他位承担权限、缓存、访问和
执行禁止等属性。因此“64 位 CPU”不能推出“可安装 \(2^{64}\) byte RAM”。
本项目的软件规格按 64 GiB 物理内存设计，架构接口保留继续扩展的能力。

\topic{64 位子模式中的分段、指令和系统入口}
在 64 位子模式中，CS、DS、ES 和 SS 的 base 按零处理，普通代码采用平坦
线性地址；FS 与 GS 仍可通过 MSR 保存 64 位 base，适合线程本地状态和
CPU-local 数据。CS 描述符的 L=1 选择 64 位子模式，且 D 必须为零。

\begin{osgridlongtable}{L{0.22\linewidth}L{0.31\linewidth}L{0.37\linewidth}}
类别 & 典型指令或形式 & 64 位子模式的重点 \\ \hline
\endfirsthead
类别 & 典型指令或形式 & 64 位子模式的重点 \\ \hline
\endhead
64 位整数 & 带 REX.W 的 MOV、ADD、SUB、IMUL、AND、OR、XOR、SHIFT &
操作 RAX--R15；许多立即数仍只有 32 位并符号扩展 \\ \hline
位置无关寻址 & RIP-relative LEA/MOV &
把下一条 RIP 加位移，便于访问附近代码和数据 \\ \hline
调用和返回 & CALL、RET、JMP、Jcc &
near 控制流使用 RIP；直接 near 分支通常仍编码相对 8/32 位位移 \\ \hline
栈 & PUSH、POP、CALL、RET &
使用 RSP，返回地址和常见栈项为 8 byte \\ \hline
快速系统调用 & SYSCALL、SYSRET；Intel 还提供 SYSENTER/SYSEXIT &
入口 MSR、权限、RCX/R11 与栈切换必须由内核安排 \\ \hline
CPU-local 切换 & SWAPGS &
在用户 GS base 与内核 GS base 之间切换，不能当作普通数据交换 \\ \hline
中断和异常 & 16-byte IDT gate、IRETQ、TSS/IST &
保存 64 位 RIP/CS/RFLAGS/RSP/SS，并可切换到专用栈 \\ \hline
扩展状态 & x87、SSE、AVX 等 &
继续由 CPUID、CR0、CR4 和 XCR0 决定可用性 \\ \hline
\end{osgridlongtable}

一些历史形式在 64 位子模式中无效，例如 PUSHA/POPA、LDS/LES、BOUND、
INTO，以及 AAA/AAS/DAA/DAS 等旧十进制调整指令。遇到具体指令时仍应查看
Intel Volume 2 或 AMD Volume 3 的 mode 表，不能只凭助记符年代猜测。

\topic{兼容子模式不是保护模式退回 32 位}
IA-32e 激活后，代码段 L=0 时运行兼容子模式。D=1 的段按 32 位默认值解码，
D=0 的段可运行 16 位代码；它们继续使用当前长模式页表。64 位操作系统因此
可以在同一分页环境下运行旧应用。

兼容子模式不能执行 64 位操作数形式或使用 R8--R15，也不能直接利用 64 位
RIP；它与普通 32 位保护模式的关键区别在于 IA-32e 已激活、分页格式已经是
长模式格式，并由 64 位系统管理全局环境。AMD 手册还明确指出，长模式不支持
real mode 或 virtual-8086 mode；要回到它们，必须先退出长模式。

本项目打开 CR0.PG 后，EFER.LMA 变为 1，但当前 CS 仍是 L=0、D=1，所以
CPU 短暂处于 32 位兼容子模式。far jump 装入 L=1 的代码段以后，目标才按
64 位子模式解码。

\topic{项目怎样从保护模式进入 64 位子模式}
项目先在 32 位保护模式完成页表、PAE 和 LME，再打开分页并远跳：

\lstinputlisting[
  language={[x86masm]Assembler},
  firstline=90,
  lastline=105,
  caption={32 位入口依次建立 PAE、LME、PG，再执行 far jump}
]{../../../../source/boot/stage1/src/entry.asm}

\lstinputlisting[
  language={[x86masm]Assembler},
  firstline=258,
  lastline=285,
  caption={PG 置位后的回读检查与 64 位入口}
]{../../../../source/boot/stage1/src/entry.asm}

\begin{osgridlongtable}{L{0.17\linewidth}L{0.28\linewidth}L{0.45\linewidth}}
状态点 & 此刻的模式与默认解码 & 为什么还没到下一状态 \\ \hline
\endfirsthead
状态点 & 此刻的模式与默认解码 & 为什么还没到下一状态 \\ \hline
\endhead
CR4.PAE=1 & 32 位保护模式 & 只允许使用 PAE/长模式页表格式，LME 仍为零 \\ \hline
CR3=PML4 & 32 位保护模式 & 页表根已选择，但 PG 尚未开启 \\ \hline
EFER.LME=1 & 32 位保护模式 & 只表示请求长模式，LMA 仍为零 \\ \hline
CR0.PG=1 & IA-32e 已激活，CS.L=0，兼容子模式 &
LMA=1，但当前代码段仍按 32 位默认值解码 \\ \hline
far jump 到 L=1 段 & 64 位子模式 &
CS 隐藏属性已经重载，目标 byte 才按 64 位环境解释 \\ \hline
\end{osgridlongtable}

顺序不能交换。LME=1 时用不合要求的分页状态打开 PG 会产生 \(\#GP\)；
页表若未映射当前 EIP、栈或 far jump 目标，开启 PG 后的取指和保存现场会
失败；把 \code{[bits 64]} 标签提前不会修复任何硬件状态。


\topic{进入长模式前要准备哪些状态}
当前实现按下面的顺序修改处理器状态：

\begin{enumerate}[label=\arabic*.]
  \item 建立 GDT，关闭中断，设置 CR0.PE，再远跳到 32 位代码段；
  \item 清零并填写初始页表；
  \item 设置 CR4.PAE，并把 PML4 物理地址写入 CR3；
  \item 通过 RDMSR/WRMSR 设置 IA32\_EFER.LME；
  \item 设置 CR0.PG，开启分页；
  \item 远跳到 GDT 中 L=1 的代码段。
\end{enumerate}

\begin{pdfdiagram}
  \node[os state] (protected) {保护模式\\CR0.PE=1\\CR0.PG=0};
  \node[os block,right=of protected] (pae) {CR4.PAE=1\\CR3=PML4};
  \node[os block,right=of pae] (lme) {IA32\_EFER\\LME=1};
  \node[os state,below=of lme] (paging) {CR0.PG=1\\LMA=1，兼容子模式};
  \node[os block,left=of paging] (jump) {far jump\\CS.L=1};
  \node[os state,left=of jump] (mode64) {64 位子模式};
  \draw[os arrow] (protected) -- (pae);
  \draw[os arrow] (pae) -- (lme);
  \draw[os arrow] (lme) -- (paging);
  \draw[os arrow] (paging) -- (jump);
  \draw[os arrow] (jump) -- (mode64);
\end{pdfdiagram}
\diagramnote{依据 Intel SDM Vol. 3A 12.8.5 重绘，并按
\filepath{source/boot/stage1/src/entry.asm} 的执行顺序标注。LME=1 时若
PAE=0 就开启分页，或者在 PG=1 时修改 LME，处理器会产生 \#GP。}

这些步骤之间没有可随意交换的空隙。分页打开后的下一次取指、栈访问和页表
读取都必须落在有效映射中。GDT 选择子错误会产生 \#GP；早期 IDT 尚未准备好
时，连续异常可能发展为三重故障，QEMU 表现为机器重新复位。

当前 Stage 1 面向项目固定的 QEMU x86-64 机器，尚未在访问 EFER 前执行
CPUID 能力检查。Kernel 入口稍后会检查处理器特性，但这无法保护更早的
RDMSR/WRMSR。若要适配未知实体 CPU，能力检查必须前移到这里。

\topic{用三张表先映射低 64 MiB}
当前 Stage 1 把 PML4、PDPT 和 PD 放在物理地址 \code{0x10000}、
\code{0x11000} 和 \code{0x12000}。PD 的一项在设置 PS 位以后映射一个
2 MiB 大页。填写前 32 项即可得到：
\[
  32\times2\ \mathrm{MiB}=64\ \mathrm{MiB}.
\]
这段恒等映射让虚拟地址和物理地址暂时相等。Stage 1、早期栈、BootInfo、
Kernel 暂存区和加载后的段都位于其中，分页开启前后的地址不会突然改变。

四级页表仍然存在，只是 PD 项直接指向大页，没有继续走到 PT。每张表含
512 个 64 位表项；地址索引分别取虚拟地址的 47:39、38:30 和 29:21 位。
“页表怎样翻译和保护地址”一章建立 4 KiB 页面时会继续使用 20:12 位的 PT 索引。

\topic{用串口定位模式转换停在哪里}
Stage 1 不在每条指令后打印日志。它只在一个状态已经写入并回读以后输出标记：

\begin{lstlisting}
[OS][STAGE1] GDT_READY
[OS][STAGE1] PROTECTED_MODE
[OS][STAGE1] PAGE_TABLES_READY
[OS][STAGE1] PAE_READY
[OS][STAGE1] LME_READY
[OS][STAGE1] PAGING_ENABLED
[OS][STAGE1] LONG_MODE
\end{lstlisting}

如果最后一行是 \code{PAE_READY}，调查范围就在 EFER、CR0.PG、当前映射和
随后的远跳之间。若连 \code{PROTECTED_MODE} 都没有出现，应先检查 GDT、
选择子和第一条 32 位入口，而不是去调试 ELF 解析。

后面几节会从磁盘字节继续往下走：先定义 Stage 1 和 Kernel 在磁盘中的位置，
再检查 ELF64 的程序头，最后把 BootInfo 地址放进 RDI 并调用 C++ 入口。
